\documentclass[12pt,a4paper,french]{book}
\usepackage[french]{babel}
\usepackage{teach}
\usepackage{array}
%%% Couleurs
\redefinecolor{thm}{title}{bgcolor}{alizarin}
\redefinecolor{prop}{title}{bgcolor}{meatbrown}
\redefinecolor{defn}{title}{bgcolor}{olivine}
\definecolor{alizarin}{rgb}{0.82, 0.1, 0.26}
\definecolor{meatbrown}{rgb}{0.9, 0.72, 0.23}
\definecolor{olivine}{rgb}{0.6, 0.73, 0.45}
%%% Leçon creuse/pleine
\usepackage{dashundergaps}
\TeacherModeOff
%\TeacherModeOn

\begin{document}

\tableofcontents

\chapter{Généralités sur les suites}


\section{Activités}

On donne dans le tableau suivant les effectifs, en million, de la population africaine depuis 1950.

\begin{center}
\begin{tabular}{|c|c|c|c|c|c|c|c|}
    \hline
    Année & 1950 & 1960 & 1970 & 1980 & 1990 & 2000 & 2010 \\
    \hline
    Population & 227,3 & 285 & 366,8 & 482,2 & 638,7 & 819,5 & 1011,2 \\
    \hline
\end{tabular}
\end{center}

\begin{act}
\begin{enumerate}
    \item Représentez cette évolution dans un repère.
    \item Calculez les coefficients multiplicateurs permettant de passer de la population d'une décennie à celle de la suivante à partir de l'année 1950.
    \item Un statisticien propose de modéliser la population africaine de la manière suivante : « À partir de 1950, tous les dix ans, la population africaine est multipliée par 1,28 ».
    \begin{enumerate}
        \item Comment justifier sa démarche ?
        \item Comparez le résultat obtenu par ce modèle en 2010 aux données réelles. Le modèle serait-il plus proche de la réalité avec un coefficient multiplicateur de 1,29 ?
    \end{enumerate}
\end{enumerate}
\end{act}

\section{Bilan et compléments}

\begin{defn}
Une \gap*{suite} numérique $u$ est une \gap*{fonction} de $\mathbb{N}$ dans $\mathbb{R}$ qui, à tout entier naturel $n$, associe son image $u(n)$, notée aussi $u_n$. $u_n$ est le \gap*{terme} de rang $n$ ; on lit « $u$ \gap*{indice} $n$ ». On note $(u_n)$ la suite.
\end{defn}

\begin{defn}
Une \gap*{suite} $(u_n)$ est définie
\begin{itemize}
    \item de manière \gap*{explicite} : s'il existe une \gap*{fonction} $f$ telle que, pour tout entier naturel $n$, $u_n = f(n)$ ;
    \item par \gap*{récurrence} : si elle est définie par la donnée de son ou de ses premiers termes et par une \gap*{relation} permettant de calculer chaque terme à partir des termes précédents.
\end{itemize}
\end{defn}

\begin{defn}
Dans un \gap*{repère} du plan, une représentation des termes de la suite $(u_n)$ est l'ensemble des \gap*{points} $M_n$ de coordonnées $(n ; u_n)$ avec $n$ entier naturel. On obtient un nuage de \gap*{points}.
\end{defn}

\begin{defn}
Une suite $(u_n)$ est
\begin{itemize}
    \item \gap*{croissante} : si, pour tout entier naturel $n$, $u_{n+1} > u_n$ ;
    \item \gap*{décroissante} : si, pour tout entier naturel $n$, $u_{n+1} \leq u_n$.
\end{itemize}
\end{defn}

\newpage

\section{Exercices}

\begin{exocor}
Une suite arithmétique $(u_n)$ est définie par $u_0 = 3$ et $u_{n+1} = u_n + 5$.

\begin{enumerate}
    \item Calculez les premiers termes de la suite $u_n$ pour $n = 0$ à $n = 4$.
    \item Déterminez la formule explicite de $u_n$ en fonction de $n$.
    \item Représentez graphiquement les termes de la suite pour $n$ de 0 à 10.
\end{enumerate}
\end{exocor}

\begin{exocor}
Une suite géométrique $(v_n)$ est définie par $v_0 = 2$ et $v_{n+1} = 2v_n$.

\begin{enumerate}
    \item Calculez les premiers termes de la suite $v_n$ pour $n = 0$ à $n = 4$.
    \item Déterminez la formule explicite de $v_n$ en fonction de $n$.
    \item Représentez graphiquement les termes de la suite pour $n$ de 0 à 10.
\end{enumerate}
\end{exocor}

\begin{exocor}
Soit une suite définie par récurrence $(w_n)$ avec $w_0 = 1$ et $w_{n+1} = w_n + n$.

\begin{enumerate}
    \item Calculez les premiers termes de la suite $w_n$ pour $n = 0$ à $n = 4$.
    \item Déterminez la formule explicite de $w_n$ en fonction de $n$.
    \item Représentez graphiquement les termes de la suite pour $n$ de 0 à 10.
\end{enumerate}
\end{exocor}

\begin{exocor}
Une suite $(x_n)$ est définie par $x_0 = 5$ et $x_{n+1} = 0.5x_n$.

\begin{enumerate}
    \item Calculez les premiers termes de la suite $x_n$ pour $n = 0$ à $n = 4$.
    \item Déterminez la formule explicite de $x_n$ en fonction de $n$.
    \item Représentez graphiquement les termes de la suite pour $n$ de 0 à 10.
\end{enumerate}
\end{exocor}

\begin{exocor}
Une suite $(y_n)$ est définie par $y_0 = 10$ et $y_{n+1} = y_n - 2$.

\begin{enumerate}
    \item Calculez les premiers termes de la suite $y_n$ pour $n = 0$ à $n = 4$.
    \item Déterminez la formule explicite de $y_n$ en fonction de $n$.
    \item Représentez graphiquement les termes de la suite pour $n$ de 0 à 10.
\end{enumerate}
\end{exocor}

\begin{exocor}
Une suite $(z_n)$ est définie par $z_0 = 4$ et $z_{n+1} = 3z_n + 1$.

\begin{enumerate}
    \item Calculez les premiers termes de la suite $z_n$ pour $n = 0$ à $n = 4$.
    \item Déterminez la formule explicite de $z_n$ en fonction de $n$.
    \item Représentez graphiquement les termes de la suite pour $n$ de 0 à 10.
\end{enumerate}
\end{exocor}

\begin{exocor}
Une suite $(a_n)$ est définie par $a_0 = 8$ et $a_{n+1} = 0.25a_n$.

\begin{enumerate}
    \item Calculez les premiers termes de la suite $a_n$ pour $n = 0$ à $n = 4$.
    \item Déterminez la formule explicite de $a_n$ en fonction de $n$.
    \item Représentez graphiquement les termes de la suite pour $n$ de 0 à 10.
\end{enumerate}
\end{exocor}

\begin{exocor}
Une suite $(b_n)$ est définie par $b_0 = 1$ et $b_{n+1} = b_n + 3n$.

\begin{enumerate}
    \item Calculez les premiers termes de la suite $b_n$ pour $n = 0$ à $n = 4$.
    \item Déterminez la formule explicite de $b_n$ en fonction de $n$.
    \item Représentez graphiquement les termes de la suite pour $n$ de 0 à 10.
\end{enumerate}
\end{exocor}

\newpage
\section{Exercices corrigés}

\begin{corrige}
\textbf{Correction de l'exercice 1 :}
\begin{enumerate}
    \item Les premiers termes de la suite $u_n$ sont $u_0 = 3$, $u_1 = 8$, $u_2 = 13$, $u_3 = 18$, $u_4 = 23$.
    \item La formule explicite de $u_n$ est $u_n = 3 + 5n$.
    \item La représentation graphique montre une droite croissante.
\end{enumerate}
\end{corrige}

\begin{corrige}
\textbf{Correction de l'exercice 2 :}
\begin{enumerate}
    \item Les premiers termes de la suite $v_n$ sont $v_0 = 2$, $v_1 = 4$, $v_2 = 8$, $v_3 = 16$, $v_4 = 32$.
    \item La formule explicite de $v_n$ est $v_n = 2^{n+1}$.
    \item La représentation graphique montre une courbe exponentielle croissante.
\end{enumerate}
\end{corrige}

\begin{corrige}
\textbf{Correction de l'exercice 3 :}
\begin{enumerate}
    \item Les premiers termes de la suite $w_n$ sont $w_0 = 1$, $w_1 = 2$, $w_2 = 4$, $w_3 = 7$, $w_4 = 11$.
    \item La formule explicite de $w_n$ est $w_n = \frac{n(n+1)}{2} + 1$.
    \item La représentation graphique montre une courbe quadratique.
\end{enumerate}
\end{corrige}

\begin{corrige}
\textbf{Correction de l'exercice 4 :}
\begin{enumerate}
    \item Les premiers termes de la suite $x_n$ sont $x_0 = 5$, $x_1 = 2.5$, $x_2 = 1.25$, $x_3 = 0.625$, $x_4 = 0.3125$.
    \item La formule explicite de $x_n$ est $x_n = 5 \times 0.5^n$.
    \item La représentation graphique montre une courbe décroissante exponentielle.
\end{enumerate}
\end{corrige}

\begin{corrige}
\textbf{Correction de l'exercice 5 :}
\begin{enumerate}
    \item Les premiers termes de la suite $y_n$ sont $y_0 = 10$, $y_1 = 8$, $y_2 = 6$, $y_3 = 4$, $y_4 = 2$.
    \item La formule explicite de $y_n$ est $y_n = 10 - 2n$.
    \item La représentation graphique montre une droite décroissante.
\end{enumerate}
\end{corrige}

\begin{corrige}
\textbf{Correction de l'exercice 6 :}
\begin{enumerate}
    \item Les premiers termes de la suite $z_n$ sont $z_0 = 4$, $z_1 = 13$, $z_2 = 40$, $z_3 = 121$, $z_4 = 364$.
    \item La formule explicite de $z_n$ est $z_n = 4 \times 3^n + \frac{3^n - 1}{2}$.
    \item La représentation graphique montre une courbe exponentielle croissante.
\end{enumerate}
\end{corrige}

\begin{corrige}
\textbf{Correction de l'exercice 7 :}
\begin{enumerate}
    \item Les premiers termes de la suite $a_n$ sont $a_0 = 8$, $a_1 = 2$, $a_2 = 0.5$, $a_3 = 0.125$, $a_4 = 0.03125$.
    \item La formule explicite de $a_n$ est $a_n = 8 \times 0.25^n$.
    \item La représentation graphique montre une courbe décroissante exponentielle.
\end{enumerate}
\end{corrige}

\begin{corrige}
\textbf{Correction de l'exercice 8 :}
\begin{enumerate}
    \item Les premiers termes de la suite $b_n$ sont $b_0 = 1$, $b_1 = 4$, $b_2 = 10$, $b_3 = 19$, $b_4 = 31$.
    \item La formule explicite de $b_n$ est $b_n = \frac{3n(n+1)}{2} + 1$.
    \item La représentation graphique montre une courbe quadratique.
\end{enumerate}
\end{corrige}

\end{document}
